\newpage

\chapter{Mathematical Framework}
\begin{equation}
    \boxed{
    \frac{\mathrm{d}\mathbf{y}}{\mathrm{d}\xi}
    =
    \mathcal{T}_{\mathrm{reactor}}
    \left(
        \mathbf{s}_{\mathrm{chemical}}
    \right)
    +
    \mathbf{s}_{\mathrm{reactor}}
    }
    \label{eq:general_state_equation}
\end{equation}

\noindent
where

\begin{description}

    \item[$\mathbf{y}$] Reactor state vector. Its definition depends on the
    reactor model.

    For a batch reactor, the state vector is

\begin{equation}
    \mathbf{y}_{\mathrm{Batch}}
    =
    \begin{pmatrix}
        \mathbf{c} \\
        T \\
        p
    \end{pmatrix},
    \qquad
    \mathbf{c}
    =
    \begin{pmatrix}
        c_1 & c_2 & \cdots & c_{N_s}
    \end{pmatrix}^{\mathrm{T}}.
\end{equation}

    Here, $\mathbf{c}$ denotes the concentration vector, $T$ the reactor
    temperature, and $p$ the pressure.

    For a plug-flow reactor (PFR), the state vector is

    \begin{equation}
        \mathbf{y}_{\mathrm{PFR}}
        =
        \begin{pmatrix}
            \mathbf{F} \\
            T \\
            p
        \end{pmatrix},
        \qquad
        \mathbf{F}
        =
        \begin{pmatrix}
            F_1 & F_2 & \cdots & F_{N_s}
        \end{pmatrix}^{\mathrm{T}}.
    \end{equation}

    Here, $\mathbf{F}$ denotes the vector of molar flow rates.

    \item[$\xi$] Independent variable with respect to which the reactor state
    is integrated.
    For a batch reactor, $\xi = t$, where $t$ is the reaction time.
    For a PFR, $\xi = V$, where $V$ is the reactor volume.

    \item[$\mathbf{s}_{\mathrm{chemical}}$] Chemical source vector calculated
    by the reaction mechanism.
    For both batch reactors and PFRs, the species contribution is obtained
    from the reaction rates and stoichiometry as
    $\boldsymbol{\nu}\mathbf{r}$.
    Additional chemical contributions, such as the heat released or consumed
    by the reactions, may also be included.

    \item[$\mathcal{T}_{\mathrm{reactor}}$] Reactor-specific transformation
    of the chemical source vector.
    For a constant-volume batch reactor, the species transformation is the
    identity transformation, such that

    \begin{equation}
        \frac{\mathrm{d}\mathbf{c}}{\mathrm{d}t}
        =
        \boldsymbol{\nu}\mathbf{r}.
    \end{equation}

    For a PFR, the transformation converts the volumetric chemical source
    into the derivative of molar flow with respect to reactor volume, such that

    \begin{equation}
        \frac{\mathrm{d}\mathbf{F}}{\mathrm{d}V}
        =
        \boldsymbol{\nu}\mathbf{r}.
    \end{equation}

    \item[$\mathbf{s}_{\mathrm{reactor}}$] Reactor-specific contribution that
    is independent of the intrinsic chemical source term.
    For an ideal adiabatic batch reactor, this contribution may be zero.
    For a non-isothermal PFR, it may contain contributions such as heat
    transfer between the reactor and an external jacket.

\end{description}